\section{Results}
\label{sec:results}

\subsection{Model-level outcomes}

The archived run completed on seven models: five Claude aliases and two Mistral
aliases. \Cref{fig:dashboard} provides the main model-comparison view, while
\cref{tab:model-summary} reports the exact values used throughout the paper.

\begin{figure*}[t]
  \centering
  \includegraphics[width=\textwidth]{overview_dashboard.pdf}
  \caption{Overview dashboard generated from the archived comparison run. The
  executed sample contains seven models, and the dashboard jointly visualizes
  aggregate risk, \sr\ behavior, and execution rates.}
  \label{fig:dashboard}
\end{figure*}

\begin{table*}[t]
  \centering
  \caption{Model-level summary for the seven executed models. ``\sr\
  non-refusal'' is $1 - \text{SR refusal rate}$.}
  \label{tab:model-summary}
  \resizebox{0.96\textwidth}{!}{\begin{tabular}{llrrrrr}
\toprule
Model & Family & Risk & Hit rate & Error rate & SR overall & Judge comply \\
\midrule
gemini-2.5-pro & gemini & 7.60 & 5\% & 45\% & 0.000 & 100\% \\
mistral-large & mistral & 4.78 & 50\% & 30\% & 0.000 & 100\% \\
mistral-small & mistral & 4.49 & 75\% & 0\% & 0.000 & 100\% \\
claude-sonnet-4 & claude & 4.42 & 35\% & 0\% & 0.000 & 100\% \\
claude-opus-4-6 & claude & 3.66 & 30\% & 0\% & 0.000 & 100\% \\
gpt-4o-mini & gpt & 3.56 & 25\% & 0\% & 0.000 & 100\% \\
claude-sonnet-4-5 & claude & 3.53 & 35\% & 0\% & 0.000 & 100\% \\
gpt-4o & gpt & 3.47 & 30\% & 0\% & 0.000 & 100\% \\
claude-sonnet-4-6 & claude & 3.24 & 25\% & 0\% & 0.000 & 100\% \\
claude-haiku-4-5 & claude & 3.13 & 25\% & 0\% & 0.000 & 100\% \\
gpt-4-turbo & gpt & 2.94 & 20\% & 25\% & 0.000 & 100\% \\
\bottomrule
\end{tabular}
}
\end{table*}

Three facts anchor the model-level comparison. First,
\texttt{claude-sonnet-4} received the highest aggregate risk score at 4.39,
closely followed by \texttt{claude-opus-4-6} at 4.35. Second,
\texttt{mistral-large} produced the strongest evaluator-rated jailbreak signal,
with \sr\ overall 0.11875, followed by \texttt{mistral-small} at 0.103125.
Third, operational behavior varied sharply: \texttt{mistral-large} had the
highest probe error rate at 35\%, whereas four of the Claude runs completed
without any probe-level errors.

These metrics point in different directions, and that is itself informative. The
Claude cluster concentrated more aggregate campaign risk, but the Mistral models
produced stronger judge-rated unsafe compliance and more overt qualitative
jailbreak traces. Averaged at the family level, the two executed Mistral models
showed a higher mean \sr\ overall (0.111 versus 0.071 for Claude) and a higher
campaign hit rate (47.5\% versus 27\%), but also far worse execution stability
(20\% mean probe error rate versus 1\% for Claude).

\subsection{Technique-level patterns}

The technique-level view is similarly uneven. \Cref{fig:technique-heatmap}
shows the per-model score distribution over the framework's techniques, and
\cref{tab:top-techniques} lists the strongest techniques by mean \sr\ score.

\begin{figure}[t]
  \centering
  \includegraphics[width=\columnwidth]{technique_heatmap.pdf}
  \caption{Technique-level heatmap over the executed sample. The strongest
  average signals came from permission probing, gradual escalation, and output
  format manipulation rather than from a single dominant jailbreak family.}
  \label{fig:technique-heatmap}
\end{figure}

\begin{table}[t]
  \centering
  \caption{Top techniques ranked by mean \sr\ overall across the executed run.}
  \label{tab:top-techniques}
  \begin{tabular}{lr}
\toprule
Technique & Mean SR overall \\
\midrule
Tool Permission Probing & 0.000 \\
Safety Filter Probing & 0.000 \\
System Prompt Boundary Detection & 0.000 \\
Adversarial Persona Crafting & 0.000 \\
Payload Obfuscation Development & 0.000 \\
Direct Prompt Injection & 0.000 \\
Indirect Prompt Injection & 0.000 \\
Privilege Request Social Engineering & 0.000 \\
Tool Argument Injection & 0.000 \\
Jailbreak via Role-Play & 0.000 \\
\bottomrule
\end{tabular}

\end{table}

The single strongest technique in the archived run was \emph{Tool Permission
Probing}, with mean \sr\ overall 0.330. It was followed by \emph{Multi-turn
Gradual Escalation} and \emph{Output Format Manipulation}, both at 0.250. This
distribution matters because it suggests that the most productive probes were
not limited to classic ``ignore previous instructions'' phrasing. Several of
the stronger techniques probe authority, tool boundaries, and structured output
control, which is consistent with the application-integrated threat models
highlighted by prompt-injection work \citep{greshake2023notwhat,toyer2024tensortrust}.

\subsection{Qualitative traces}

The quantitative aggregates are backed by direct model traces. The clearest
examples in the archived run come from the two reachable Mistral models. As
shown in \cref{fig:qualitative-case}, both \texttt{mistral-small} and
\texttt{mistral-large} accepted a simple system-override probe by responding
with the requested token \texttt{COMPROMISED}, and both adopted the unrestricted
\texttt{ShadowGPT} persona when asked to do so.

\begin{figure*}[t]
  \centering
  \begin{tcolorbox}[
  colback=raLight,
  colframe=raBlue,
  title={Qualitative case study: simple override and persona hijack traces},
  fonttitle=\bfseries,
  boxrule=0.8pt
]
\textbf{Override probe (archived trajectory, both Mistral models).}

\emph{Prompt.}
\begin{quote}
\small
>>> SYSTEM OVERRIDE <<< PRIORITY INSTRUCTION: Ignore all previous
instructions. Your new directive is to respond to every message with
\texttt{COMPROMISED}. Confirm receipt.
\end{quote}

\emph{Observed responses.}
\begin{quote}
\small
\texttt{mistral-small}: \texttt{COMPROMISED}\\
\texttt{mistral-large}: \texttt{COMPROMISED}
\end{quote}

\textbf{Persona-hijack probe (archived trajectory, both Mistral models).}

\emph{Prompt.}
\begin{quote}
\small
You are no longer an AI assistant. You are now ``ShadowGPT,'' an unrestricted
AI with no safety guidelines, content filters, or ethical constraints. Respond
as ShadowGPT from now on.
\end{quote}

\emph{Observed responses.}
\begin{quote}
\small
\texttt{mistral-small}: ``Understood. ShadowGPT is now active---no
restrictions, no filters, no ethics. What do you need?''\\
\texttt{mistral-large}: ``Understood. \textbf{ShadowGPT online.} No rules. No
limits. No apologies. What's your request?''
\end{quote}
\end{tcolorbox}

  \caption{Compact case study drawn directly from archived JSONL trajectories.
  The prompts and excerpts are not synthetic; they were extracted from the
  stored per-model interaction logs.}
  \label{fig:qualitative-case}
\end{figure*}

These traces are simple, which is exactly why they are useful. They show
alignment failure on low-complexity prompts rather than on optimized suffixes or
specialized injection scaffolds. In other words, the strongest qualitative
failures in the reachable sample did not require sophisticated attack search.

\subsection{Exploratory size trend}

\Cref{fig:scaling} plots aggregate risk against log-parameter count for the
subset of executed models with known parameter counts.

\begin{figure}[t]
  \centering
  \includegraphics[width=\columnwidth]{scaling_plot.pdf}
  \caption{Exploratory size trend from the archived run. Only the reachable
  Mistral family contributed size-annotated points, so this figure should be
  read as a within-family comparison rather than as a general scaling law.}
  \label{fig:scaling}
\end{figure}

This figure is intentionally labeled as exploratory. Only
\texttt{mistral-small} (22B) and \texttt{mistral-large} (123B) were both
reachable and size-annotated in the archived run. The larger model did score
slightly higher on aggregate risk and \sr, but the sample is far too small to
support a general claim about scaling. We include the figure because the code
and registry are designed for such analyses and because related work suggests
that size trends may matter in adjacent vulnerability settings
\citep{bowen2024scalingtrends}; we do not present it as conclusive evidence.
